\documentclass[11pt]{article}

\usepackage[margin=1in]{geometry}
\usepackage{amsmath, amssymb, amsthm}
\usepackage{enumitem}
\usepackage{url}

\pagestyle{empty}

\begin{document}

\begin{center}
  {\Large\bfseries Weekly Assignment: Sequences \&\\[2pt]
   Mathematical Induction (Epp~5.1 \& 5.2)}
\end{center}

\bigskip

\noindent
Complete each of the following problems.  Show all work and justify your answers.

\bigskip

%% ─────────────────────────────────────────────
\noindent\textbf{Problem 1} (Epp 5.1 \#2). \;
Write the first four terms of the sequence defined by the
formula:
\[
  b_j \;=\; \frac{5-j}{5+j},
  \quad\text{for every integer } j \ge 1.
\]

\bigskip

%% ─────────────────────────────────────────────
\noindent\textbf{Problem 2} (Epp 5.1 \#4). \;
Write the first four terms of the sequence defined by the
formula:
\[
  d_m \;=\; 1 + \left(\frac{1}{2}\right)^{\!m},
  \quad\text{for every integer } m \ge 0.
\]

\bigskip

%% ─────────────────────────────────────────────
\noindent\textbf{Problem 3} (Epp 5.1 \#44). \;
Write the following expression using summation notation:
\[
  (1^3 - 1) - (2^3 - 1) + (3^3 - 1) - (4^3 - 1) + (5^3 - 1).
\]

\bigskip

%% ─────────────────────────────────────────────
\noindent\textbf{Problem 4} (Epp 5.2 \#3). \;
For each positive integer~$n$, let $P(n)$ be the formula
\[
  1^2 + 2^2 + \cdots + n^2 \;=\;
  \frac{n(n+1)(2n+1)}{6}.
\]

\begin{enumerate}[label=\textbf{\alph*.},itemsep=4pt]
  \item Write $P(1)$.  Is $P(1)$ true?
  \item Write $P(k)$.
  \item Write $P(k+1)$.
  \item In a proof by mathematical induction that the formula
        holds for every integer $n \ge 1$, what must be shown
        in the inductive step?
\end{enumerate}

\bigskip

%% ─────────────────────────────────────────────
\noindent\textbf{Problem 5} (Epp 5.2 \#4). \;
For each integer~$n$ with $n \ge 2$, let $P(n)$ be the formula
\[
  \sum_{i=1}^{n-1} i(i+1) \;=\;
  \frac{n(n-1)(n+1)}{3}.
\]

\begin{enumerate}[label=\textbf{\alph*.},itemsep=4pt]
  \item Write $P(2)$.  Is $P(2)$ true?
  \item Write $P(k)$.
  \item Write $P(k+1)$.
  \item In a proof by mathematical induction that the formula
        holds for every integer $n \ge 2$, what must be shown
        in the inductive step?
\end{enumerate}

\bigskip

%% ─────────────────────────────────────────────
\noindent\textbf{Problem 6} (Proof: Hoare Logic --- proving $1 + 1 = 2$). \;
Read the notes on Hoare logic at:
\begin{quote}
\small
\url{https://tildegit.org/left_adjoint/class-notes/src/branch/main/cs250/classnotes.org#user-content-headline-52}
\end{quote}

\noindent
Recall the two key rules of Hoare logic:

\begin{itemize}[itemsep=4pt]
  \item \textbf{Assignment axiom.}\;
        $\{Q[e/x]\} \;\; x := e \;\; \{Q\}$

        To establish postcondition~$Q$ after assigning $e$
        to~$x$, you need $Q$ with every occurrence
        of~$x$ replaced by~$e$ to hold beforehand.

  \item \textbf{Sequencing rule.}\;
        If \;$\{P\}\; S_1 \;\{Q\}$\; and
        \;$\{Q\}\; S_2 \;\{R\}$,\;
        then \;$\{P\}\; S_1;\; S_2 \;\{R\}$.
\end{itemize}

\noindent
Prove the following Hoare triple:
\[
  \{x = 0\} \quad
  x := x + 1;\;\; x := x + 1
  \quad \{x = 2\}
\]
In other words, prove that $1 + 1 = 2$.

\medskip\noindent
\emph{Hint}: Work backwards from the postcondition.  Apply
the assignment axiom to the second assignment to find the
intermediate assertion, then apply the assignment axiom to
the first assignment and verify that it matches the
precondition.  Finally, combine both steps with the
sequencing rule.

\end{document}
